\section{PLANO DETALHADO DE VALIDAÇÃO}
\label{ap:validacao}

\begin{longtable}{p{0.08\textwidth}p{0.19\textwidth}p{0.08\textwidth}p{0.28\textwidth}p{0.27\textwidth}}
\caption{Casos de teste funcional do SIGE Desk}\label{tab:casos-teste}\\
\hline
\textbf{ID} & \textbf{Requisitos} & \textbf{Prior.} & \textbf{Cenário} & \textbf{Resultado esperado} \\
\hline
\endfirsthead
\hline
\textbf{ID} & \textbf{Requisitos} & \textbf{Prior.} & \textbf{Cenário} & \textbf{Resultado esperado} \\
\hline
\endhead
CT-01 & RF-05, RF-06, RF-21 & Alta & Cliente abre ticket de alteração de campanha, inclusive sem campanha cadastrada. & Ticket recebe identificador, status Aberta, urgência, prazo desejado e SLA pendente de classificação; campanha pendente segue para triagem. \\
CT-02 & RF-07, RF-09, RN-03, RN-14 & Alta & Atendimento inicia triagem, classifica e atribui o ticket. & Prioridade, regra e prazos de SLA calculados desde a abertura e responsável ficam registrados; confirmação automática não conta como primeira resposta. \\
CT-03 & RF-08, RF-09, RF-10 & Alta & Gestor atualiza o status e inclui comentário. & Histórico registra usuário, data, hora e alteração realizada. \\
CT-04 & RF-11, RN-08 & Média & Ticket aguarda complemento na triagem e na execução; Cliente envia informação. & Motivo, origem e pausa ficam visíveis; somente complemento suficiente confirmado retoma a fase correta, sem permitir cancelamento após execução. \\
CT-05 & RF-11, RF-12, RN-04, RN-05, RN-06, RN-11 & Alta & Gestor registra entrega de tipo com e sem aprovação; Cliente corrige e reabre concluída. & Tipo obrigatório entra em validação e só Cliente aprova/corrige; tipo sem aprovação conclui com dispensa automática; correção e reabertura preservam histórico e ciclos corretos. \\
CT-06 & RF-14, RF-16, RF-21, RN-12 & Média & Usuário filtra demandas por abertura e situação de SLA. & Listagem distingue prazo pendente, pausa, vencimento de primeira resposta e vencimento de resolução, no escopo autorizado. \\
CT-07 & RF-20, RN-13 & Alta & Ticket é atribuído, comentado, colocado em aguardo, tem complemento confirmado, entra em validação, vence, conclui, reabre ou é cancelado. & Envolvidos autorizados recebem notificação no sistema sem revelação de conteúdo fora da permissão. \\
CT-08 & RF-02, RN-09, RNF-02 & Alta & Cliente tenta consultar ticket de outra organização. & Acesso é negado e a tentativa é registrada conforme a política de auditoria. \\
CT-09 & RF-22, RN-05, RN-10 & Alta & Administrador configura um tipo de demanda, inclusive Aprovação de criativo. & Campos tipados, regras de aprovação/evidência e material submetido são aplicados a novos tickets sem exigir decisão antes da validação. \\
CT-10 & RF-23 & Alta & Administrador configura calendário e prazos do SLA. & Horário, fuso, feriados, recessos, pausa em validação e precedência por cliente/tipo são usados em novos tickets; os existentes preservam a regra registrada. \\
CT-11 & RF-01, RF-02 & Alta & Usuário realiza autenticação e acessa o sistema com seu perfil. & Usuário visualiza somente funções e tickets autorizados. \\
CT-12 & RF-03, RF-04 & Alta & Administrador/Atendimento cadastra campanha e administra inativação. & Campanha fica vinculada ao cliente; cliente/campanha inativos não recebem novas solicitações e inativações que deixam pendências sem responsável/Cliente são bloqueadas. \\
CT-13 & RF-09, RF-17, RNF-04 & Alta & Usuário consulta o histórico após mudanças de status, responsável, prioridade, prazo e aprovação. & Histórico exibe autor, data, hora, motivo, campo, valor anterior e novo valor. \\
CT-14 & RF-13, RN-07 & Média & Atendimento cancela ticket antes da execução. & Motivo do cancelamento fica registrado e o ticket não segue para execução. \\
CT-15 & RF-15 & Média & Usuário autorizado consulta o painel de demandas. & Painel apresenta totais por status, prioridade, responsável e prazo. \\
CT-16 & RF-18 & Baixa & Solicitante informa métricas de contexto em ticket aplicável. & Métricas informadas ficam registradas no ticket. \\
CT-17 & RF-19 & Baixa & Usuário exporta uma listagem filtrada. & Exportação contém somente tickets retornados pelos filtros e autorizados ao perfil. \\
\hline
\multicolumn{5}{l}{\footnotesize\textbf{Fonte:} Elaborado pelo grupo a partir do plano de testes preliminar.}
\end{longtable}

Cada execução deverá registrar pré-condição, passos, resultado obtido, evidências e situação final. O conjunto será considerado aprovado quando todos os casos de alta prioridade forem aprovados, pelo menos 90\% dos dezessete casos forem aprovados e não houver defeito crítico em aberto. Casos reprovados deverão gerar correção ou justificativa documentada de adiamento.
